\documentclass[11pt,a4paper,twoside,pdf]{article}

% Paquetes (añade otros si los necesitas):

\usepackage[T1]{fontenc}
\usepackage[utf8]{inputenc}


\usepackage[pdftex,final]{graphicx}
\bibliographystyle{plain} % We choose the "plain" reference style

% Style package
% Font Package (Palatino)
\usepackage{mathpazo}
% Packages for specific capabilities
\usepackage{rotating} % for text rotation in tables
\usepackage{multirow} % for multirow in tables
\usepackage{subfigure} % Place subfigures in figure environment
% Packages for specific symbols
\usepackage{amssymb}
\usepackage{amsmath}
\usepackage{amsfonts}
\usepackage{braket}
\usepackage{eurosym} % Euro symbol
\usepackage{bbding} % for \XSolidBrush
\usepackage{pifont} % for \ding{55} (a check mark)
\usepackage{macro}


\usepackage{latexsym}
\usepackage{soul}
\usepackage{array}
%\usepackage{marvosym}
\usepackage{epsfig}
\usepackage{graphics}
\usepackage{amsfonts}
\usepackage{xspace}
\usepackage{color}
\usepackage{booktabs}
\usepackage{xtab}
\usepackage[colorlinks=true,urlcolor=blue,linkcolor=blue,citecolor=blue]{hyperref}
\numberwithin{equation}{section}


\linespread{1.05}

% TFG en inglés:
\usepackage[english]{babel} 
\addto\captionsenglish{\renewcommand{\chaptername}{}}

% TFG en español:
%\usepackage[spanish,es-nodecimaldot,es-tabla,es-lcroman,es-nosectiondot,es-noindentfirst]{babel}
%\renewcommand\spanishchaptername{}

% Formato de la página:
\usepackage{fancyhdr}
\usepackage[top=2.88cm,bottom=2.97cm,left=2.95cm,right=2.95cm]{geometry}
\setlength{\parskip}{0.1cm}

% Pon aquí tus definiciones:

\newcommand{\dis}{\displaystyle}
%\sodef\an{}{.2em}{1em plus1em}{2em plus.1em minus.1em}

\begin{document}

% Portada %%%%%%%%%%%%%%%%%%%%%%%%%%%%%%%%%%%%%%%%%%%%%%%%%%%%%%%%%%%%%%%%%%%%%%

\pagestyle{empty}


\noindent
\begin{tabular}{r}
\includegraphics[width=8.8cm]{escudoUGRmonocromo.png} \\[-1.8ex]
\hspace{31mm}\vspace{-8mm}
\begin{tabular}{c}
\hline\\[-1ex]\hskip-2mm
{\bf Facultad de Ciencias}\hspace{18mm}
\end{tabular}
\end{tabular}

\large
\vspace{30mm}
\hspace{25mm}
\begin{tabular}{l}
GRADO EN F\'ISICA
\end{tabular}

\vspace{45mm}
\hspace{25mm}
\begin{tabular}{l}
TRABAJO FIN DE GRADO
\\[1.5ex]
\LARGE\bf Herramientas para cálculos \\ 
\LARGE\bf perturbativos en renormalización \\
\LARGE\bf de Hamiltonianos \\
\end{tabular}
%
\vfill
\hspace{25mm}
\begin{tabular}{l}
Presentado por:
\\
{\bf D. ZhuoZhuo Liu}
\\[3ex]
Curso Académico 2024/2025
\end{tabular}
%

\newpage
%
\begin{center}
{\bf Resumen}
\bigskip

\begin{minipage}{0.8\linewidth}
bla bla bla bla bla bla bla bla bla bla bla bla bla bla bla
bla bla bla bla bla bla bla bla bla bla bla bla bla bla bla
bla bla bla bla bla bla bla bla bla bla bla bla bla bla bla 
bla bla bla bla bla bla bla bla bla bla bla bla bla bla bla
bla bla bla bla bla bla bla bla bla bla bla bla bla bla bla
\end{minipage}

\newpage

{\bf Abstract} 
\bigskip

\begin{minipage}{0.8\linewidth}
bla bla bla bla bla bla bla bla bla bla bla bla bla bla bla
bla bla bla bla bla bla bla bla bla bla bla bla bla bla bla
bla bla bla bla bla bla bla bla bla bla bla bla bla bla bla 
bla bla bla bla bla bla bla bla bla bla bla bla bla bla bla
bla bla bla bla bla bla bla bla bla bla bla bla bla bla bla
\end{minipage}

\newpage

\end{center}

% Indice %%%%%%%%%%%%%%%%%%%%%%%%%%%%%%%%%%%%%%%%%%%%%%%%%%%%%%%%%%%%%%%%%%%%%%%
%\newpage

\tableofcontents

% Texto %%%%%%%%%%%%%%%%%%%%%%%%%%%%%%%%%%%%%%%%%%%%%%%%%%%%%%%%%%%%%%%%%%%%%%%%
%\newpage

\pagestyle{fancy}
\fancyhead[RO,LE]{\leftmark}
\fancyhead[LO,RE]{\thepage}
\fancyfoot{}

\newpage

\section{Introduction}



It's in our interest to formulate a theory compatible with special relativity and
quantum mechanics, this culminates in quantum field theories. 

we should expect the most general theories to be scale invariant, 

In the field of theoretical physics, divergences are a common occurrence, 

In the context of theoretical physics renormalization is an important procedure 

Tipical renormalization procedures apply to Lagrangian dynamics, 

\newpage

\section{Theoretical background}

To describe the relativistic interactions of effective particles in the framework of 
the renormalization group procedure for effective particles, making use of quantum 
theory in the front form of Hamiltonian dynamics. 

\subsection{Canonical Hamiltonian}\label{sec:canonical_hamiltonian}

\subsubsection{Classical mechanics}

The canonical Hamiltonian is a function of the canonical coordinates $q_i$,
the canonical momenta $p_i$, and time $t$. The Hamiltonian is a function that
describes the total energy of the system, and is obtained from applying the Legendre 
transformation to the Lagrangian of the system, this is defined as,
\begin{equation}
    H(q,p,t) = \sum_i p_i \dot{q}_i - L(q,\dot{q},t),
\end{equation}
where $H$ is the Hamiltonian, $L$ is the Lagrangian, $q_i$ are the canonical
coordinates, $\dot{q}_i$ are the canonical velocities, and $p_i$ are the
canonical momenta. The canonical momenta are defined as,
\begin{equation}
    p_i = \frac{\partial L}{\partial \dot{q}_i},
\end{equation}

The canonical Hamiltonian governs the time evolution of the system via the 
Hamilton equations of motion, which are given by,
\begin{equation}
    \dot{q}_i = \frac{\partial H}{\partial p_i}, \quad \dot{p}_i = -\frac{\partial H}{\partial q_i}.
\end{equation}

\subsubsection{Field theory}

In the case of field theories, a Hamiltonian density can be defined such that 
integrating over space coordinates the Hamiltonian is obtained. 

\begin{equation}
    H = \int d^3x \mathcal{H},
\end{equation}
for simplicity we will call the Hamiltonian density $\mathcal{H}$, as Hamiltonian.

For a classical field $\phi(x)$, with the Lagrangian density $\mathcal{L}
(\phi,\partial_\mu\phi)$, the Hamiltonian density is defined as,

\begin{equation}
    \mathcal{H} = \pi \dot{\phi} - \mathcal{L},
\end{equation}
where $\pi$ is the canonical momentum, defined as,

\begin{equation}
    \pi = \frac{\partial \mathcal{L}}{\partial \dot{\phi}}.
\end{equation}

\subsubsection{Quantum field theory}

In quantum mechanics, or quantum field theory, the Hamiltonian becomes an operator 
acting on a Hilbert space. 

Solving the system means obtaining the spectrum of the theory, this is finding the 
eigenenergies of the system. In general this is a non trivial process. 

In field theory calculations the fundamental quantities to consider are the 
energy-momenta, angular momentum, and boosts, due to its symmetries
under the Poincaré group.

The Poincaré group contains the Lorentz group and the translations, the generators of
the Poincaré group are the four-momentum operator $P^\mu$ and the angular momentum
$M^{\mu\nu}$, which satisfy the following commutation relations,
\begin{align}
    \sbrackets{P^\mu, P^\nu} &= 0, \\
    \sbrackets{M^{\mu\nu}, P^\rho} &= 
    i\left( g^{\nu\rho}P^\mu - g^{\mu\rho}P^\nu \right), \\
    \sbrackets{M^{\mu\nu}, M^{\rho\sigma}} &    = i\left( g^{\nu\rho}M^{\mu\sigma} 
    - g^{\mu\rho}M^{\nu\sigma} + g^{\sigma\mu}M^{\nu\rho} - g^{\sigma\nu}
    M^{\mu\rho}\right).
\end{align}


\subsection{Fock space}

In order to describe a variable number of particles in our system, the use of the Fock
space is needed. Intoduced by V.A. Fock in 1932 \cite{1932ZPhy...75..622F}, the Fock 
space is a sum of a set of Hilbert spaces, each one corresponding to a different number of
particles in the system. 

The Fock space is defined as the direct sum of tensor products of the single 
particle Hilbert space $\mathbb{H}$,

\begin{equation}
    \mathbb{F}_\nu = \bigoplus_{n=0}^{\infty} S_\nu \mathbb{H}^{\otimes n} = 
    \mathbb{C} \oplus \mathbb{H} \oplus S_\nu(\mathbb{H} \otimes \mathbb{H}) 
    \oplus S_\nu(\mathbb{H} \otimes \mathbb{H}\otimes \mathbb{H}) \oplus \cdots,
\end{equation}
where $S_\nu$ is the symmetrization operator depending on whether the particles described
are bosonic or fermionic, it symmetrizes or antisymmetric the tensors, and 
$\mathbb{C}$ is the complex scalar, corresponding to the states with no particles. 

This way a general state in the Fock space can be expressed as,

\begin{equation}
    \ket{\Psi}_\nu = \ket{\Psi_0}_\nu \oplus \ket{\Psi_1}_\nu \oplus
    \ket{\Psi_2}_\nu \oplus \cdots = a \ket{0}_\nu + \sum_{i=1}a_i \ket{\psi_i}_\nu +
    \sum_{i,j=1}a_{ij} \ket{\psi_i \psi_j}_\nu + \cdots,
\end{equation}
where $\ket{\Psi_0}_\nu$ is the vacuum state, $\ket{\Psi_1}_\nu$ is the one particle
state, $\ket{\Psi_2}_\nu$ is the two particle state, and so on. The coefficients $a_i$
are the amplitudes of the states, in general complex numbers. 

In the case of QCD, a geneal state in the Fock space involves a superposition of all 
possible multiparticle states, build from quarks, antiquarks, and gluons, with the
correct quantum numbers. This way a general state in the Fock space can be expressed as,
\begin{equation}
    \ket{\Psi} = c_0 \ket{0} + c_1 \ket{q\bar{q}} + c_2 \ket{q\bar{q}g} + c_3 \ket{qqq} 
    + c_4 \ket{gg} + \cdots,
\end{equation}

\subsection{Front form of Hamiltonian Dynamics}

The front form of dynamics introduced by Dirac (1949) \cite{dirac_front_forms_1949} 
offers a couple advantages to the 

The quantization hypersurface considered is,

\begin{equation}
    x^+ = t+z = x^0 + x^3 = 0,
\end{equation}
then the rest of coordinates will be defined as

\begin{equation}
    x^- = x^0-x^3, \quad x^\perp = (x^1, x^2),
\end{equation}

In this set of coordinates, the fundamental quantities 

The Hamiltonian in FF quantization is obtained from the Lagrangian density using 
the Legendre transformation, similar to the procedure explained in the section
\ref{sec:canonical_hamiltonian}, but with respect to the new coordinates.

\subsection{Counterterms}

The counterterms is an additional term added to the initial or bare Hamiltonian 
$\mathcal{H}_0$ to deal with the divergences due to loops, produced during the 
process. This way, ensure that the effective Hamiltonian $\mathcal{H}_t$ remains 
finite at all values of $t$.

\subsection{RGPEP}
The renormalization group procedure for effective particles (RGPEP), as it's name 
indicate is a renormalization group procedure applied to the Hamiltonian formulation.

By considering a series of unitary transformations, the RGPEP is able to construct a 
series of effective Hamiltonians $\mathcal{H}_s$, and the corresponding effective 
particles, by the use of effective particle operators (namely creation and 
annihilation operators) that differs from the canonical ones by the unitary 
transformation $\mathcal{U}_s$,

\begin{equation}
    a_s = \mathcal{U}_sa_0\mathcal{U}_s^\dagger,
    \label{eq:effective_particle_operator}
\end{equation}
labeled by the parameter $s$, the renormalization group parameter. This parameter 
has dimension of length. Physically $s$ has the interpretation of the characteristic 
size of the effective particles.

Due to dimensional and notational reasons that will be explained latter, it's 
convenient to consider the scale parameter $t = s^4$ instead of $s$. 

The effective Hamiltonian is related to the regulated canonical one with counter-terms 
by the condition,

\begin{equation}
    \mathcal{H}_t(a_t) = \mathcal{H}_0(a_0)
\end{equation}



Then $s=0$ correspond to the point-like or bare particles, and recovering the original

combining with the equation \eqref{eq:effective_particle_operator}, and considering 
the parameter $t$ instead of $s$, the condition becomes,

\begin{equation}
    \mathcal{H}_t(a_0) = \mathcal{U}_t^\dagger\mathcal{H}_0(a_0) \mathcal{U}_t,
\end{equation}
differentiating with respect of $t$, 

\begin{equation}
    \mathcal{H}_t^\prime (a_0) \equiv \frac{d}{dt}\mathcal{H}_t(a_0) = 
    \sbrackets{-\mathcal{U}_t^\dagger \mathcal{U}_t^\prime, \mathcal{H}_t(a_0)} 
    = \sbrackets{\mathcal{G}_t(a_0), \mathcal{H}_t(a_0)},
\end{equation}
where $\mathcal{G}_t$ is the RGPEP generator, this generator is defined as 

We will consider the generator from Ref. \cite{PEP}

\begin{equation}
    \mathcal{G}_t = \sbrackets{\mathcal{H}_f, \mathcal{H}_{Pt}},
\end{equation}
where $\mathcal{H}_f$, the free part of $\mathcal{H}_t$, and $\mathcal{H}_{Pt}$ is 
defined as function of the interacting term.

\begin{equation}
    =  \sbrackets{\sbrackets{\mathcal{H}_f, \mathcal{H}_{Pt}}, \mathcal{H}_t}
\end{equation}

where $\mathcal{H}_t$ is the Hamiltonian interested in solving.

The free Hamiltonian $\mathcal{H}_f$ is the part of $\mathcal{H}_0(a_0)$ that does 
not depend on the coupling constants, 


\subsection{Diagram representation}

From the expression of the Hamiltonian, different terms can be separated and correlate
to a diagram representation of the process, similar to the Feynman diagrams.




\section{Cases}

\subsection{Scalar case}


\section{Code implementation}

\section{Diagrams obtained}

\subsection{Scalar}

\subsection{QED}

\subsection{QCD}

\section{Conclusions and future work}



% Referencias %%%%%%%%%%%%%%%%%%%%%%%%%%%%%%%%%%%%%%%%%%%%%%%%%%%%%%%%%%%%%%%%%
%\newpage

\bibliography{references} % Entries are in the references.bib file
\nocite{QFT}
\nocite{QCDG}

\end{document}

